\documentclass[11pt,a4paper]{article}

\usepackage[utf8]{inputenc}
\usepackage[T1]{fontenc}
\usepackage{XCharter}
\usepackage[brazil]{babel}
\usepackage[margin=2cm]{geometry}
\usepackage{amsmath, amssymb}
\usepackage[xcharter,bigdelims,vvarbb]{newtxmath}
% newtxmath's operators font requests the fontspec family `XCharter(0)' in OT1;
% without these substitutions math digits and operator names silently fall back
% to Computer Modern (LaTeX Font Warning: Font shape `OT1/XCharter(0)/m/n' undefined).
\DeclareFontFamily{OT1}{XCharter(0)}{}
\DeclareFontShape{OT1}{XCharter(0)}{m}{n}{<->ssub * XCharter-TLF/m/n}{}
\DeclareFontShape{OT1}{XCharter(0)}{m}{it}{<->ssub * XCharter-TLF/m/it}{}
\DeclareFontShape{OT1}{XCharter(0)}{b}{n}{<->ssub * XCharter-TLF/b/n}{}
\DeclareFontShape{OT1}{XCharter(0)}{bx}{n}{<->ssub * XCharter-TLF/b/n}{}
\usepackage{enumerate}
\usepackage{xcolor}

\pagestyle{empty}
\setlength{\parindent}{0pt}
\setlength{\parskip}{0.5em}

\newcommand{\thetav}{\boldsymbol{\theta}}

\begin{document}

%% =============================================================
%%  Header
%% =============================================================
{\Large\bfseries Aprendizado Profundo \textemdash{} Gabarito} \hfill \textbf{Prova, Unidades 4 e 5}\\
Prof. Lucas S. Kupssinskü

\rule{\linewidth}{0.8pt}

%% =============================================================
%%  Objective questions
%% =============================================================
\section*{Questões Objetivas}

\textbf{Quadro de respostas:}
\begin{center}
\begin{tabular}{|c|c|c|c|c|c|c|c|c|c|}
\hline
1 & 2 & 3 & 4 & 5 & 6 & 7 & 8 & 9 & 10 \\\hline
f & c & a & g & d & b & h & e & c & a \\\hline
\end{tabular}
\end{center}

\color{blue}

%% ====================== UNIDADE 4 ======================

\subsection*{Questão 1 \textemdash{} Resposta: (f)}
Com $H = 8$ e $d_{\text{model}} = 512$, cada cabeça tem $d_k = 512/8 = 64$. Logo $\mathbf{Q}_i = \mathbf{X} \mathbf{W}_Q^i \in \mathbb{R}^{10 \times 64}$; os \textit{scores} $\mathbf{Q}_i \mathbf{K}_i^\top \in \mathbb{R}^{10 \times 10}$ (uma entrada por par de \textit{tokens}); após concatenar as $8$ cabeças ($10 \times 512$) e aplicar $\mathbf{W}_O$ ($512\times 512$), a saída é $10 \times 512$.\\
\textbf{(a)} usa $\mathbf{Q}_i$ com $d_{\text{model}}$ inteiro (não dividiu por $H$).
\textbf{(b)} \textit{scores} $64\times 64$ confunde dimensão de chave com comprimento de sequência.
\textbf{(c)} saída $10\times 64$ esquece a concatenação das cabeças.
\textbf{(d)/(e)/(g)} trocam dimensões de sequência, cabeças e chave.
\textbf{(h)} saída $512\times 512$ confunde a matriz $\mathbf{W}_O$ com a saída.

\subsection*{Questão 2 \textemdash{} Resposta: (c)}
Cada projeção de atenção é $64\times 64 + 64 = 4160$; com $\mathbf{W}_Q, \mathbf{W}_K, \mathbf{W}_V, \mathbf{W}_O$: $4 \times 4160 = 16\,640$ (o número de cabeças $H$ \emph{não} altera o total, pois as cabeças particionam $d_{\text{model}}$). FFN: $\mathbf{W}_1 = 64\times 256 + 256 = 16\,640$ e $\mathbf{W}_2 = 256\times 64 + 64 = 16\,448$, total $33\,088$. Dois \textit{LayerNorm}: $2 \times (64 + 64) = 256$. Soma: $16\,640 + 33\,088 + 256 = 49\,984$.\\
\textbf{(a)} $49\,408$ ignora os vieses das projeções e da FFN (mantém os \textit{LayerNorm}).
\textbf{(b)} $45\,824$ esquece a projeção de saída $\mathbf{W}_O$.
\textbf{(d)} $49\,152$ ignora todos os vieses e os \textit{LayerNorm}.
\textbf{(e)} $33\,088$ conta apenas a FFN.
\textbf{(f)} $16\,640$ conta apenas a atenção.
\textbf{(g)} $49\,728$ esquece os \textit{LayerNorm}.
\textbf{(h)} $50\,240$ conta os \textit{LayerNorm} em dobro.

\subsection*{Questão 3 \textemdash{} Resposta: (a)}
\textbf{I (V):} as componentes do produto escalar somam variância $d$; dividir por $\sqrt{d}$ devolve variância $\approx 1$, evitando que o \textit{softmax} sature.
\textbf{II (V):} sem as projeções $\mathbf{W}_Q, \mathbf{W}_K, \mathbf{W}_V, \mathbf{W}_O$, a atenção ``crua'' é só \textit{softmax} e somas ponderadas, sem parâmetros.
\textbf{III (V):} \textit{self-attention} é permutação-equivariante; sem \textit{positional encoding} a ordem dos \textit{tokens} não seria distinguida.
\textbf{IV (F):} dividir por $d$ levaria a variância a $d/d^2 = 1/d$ (\textit{softmax} ``frio'' demais); o correto é $\sqrt{d}$.
\textbf{V (F):} a máscara causal soma $-\infty$ às posições futuras \emph{antes} do \textit{softmax}; o peso $0$ aparece \emph{depois} do \textit{softmax}.
Verdadeiras: I, II e III.\\
\textbf{(b)/(d)} incluem IV (falsa).
\textbf{(e)/(f)} incluem V (falsa).
\textbf{(c)} omite II.
\textbf{(g)} marca tudo.
\textbf{(h)} omite III.

\subsection*{Questão 4 \textemdash{} Resposta: (g)}
\textbf{I (V):} o termo $\mathbf{Q}\mathbf{K}^\top$ é $O(n^2 d)$; dobrar $n$ quadruplica a parcela $n^2$.
\textbf{II (V):} com \textit{KV-cache}, cada novo \textit{token} reaproveita $\mathbf{K}, \mathbf{V}$ do prefixo, custando $O(n)$, mas a memória do \textit{cache} cresce $\propto n$.
\textbf{III (V):} \textit{cross-attention} usa \textit{query} do \textit{decoder} e \textit{key}/\textit{value} do \textit{encoder}.
\textbf{IV (V):} o \textit{positional encoding} sinusoidal é uma função fixa de senos/cossenos, sem parâmetros treináveis.
\textbf{V (F):} as projeções de atenção totalizam $\approx 4\,d_{\text{model}}^2$ \emph{independentemente} de $H$, pois as cabeças particionam $d_{\text{model}}$; aumentar $H$ não aumenta o total.
Verdadeiras: I, II, III e IV.\\
\textbf{(a)/(b)/(d)/(f)} omitem alguma afirmação verdadeira.
\textbf{(c)/(e)/(h)} incluem V (falsa).

\subsection*{Questão 5 \textemdash{} Resposta: (d)}
Aplicando $\min(1,\, p/q)$ ao \textit{token} proposto em cada posição:
\begin{itemize}
    \item Pos.\ 1: $p = 0.90 \ge q = 0.60 \Rightarrow$ aceito com prob.\ $1$ (sempre).
    \item Pos.\ 2: $p = 0.70 \ge q = 0.50 \Rightarrow$ sempre aceito.
    \item Pos.\ 3: $p = 0.20 < q = 0.80 \Rightarrow$ aceito com prob.\ $0.20/0.80 = 0.25$ (rejeitado com prob.\ $0.75$).
    \item Pos.\ 4: $p = 0.60 \ge q = 0.40 \Rightarrow$ sempre aceito (se alcançada).
    \item Pos.\ 5: $p = 0.35 < q = 0.70 \Rightarrow$ aceito com prob.\ $0.5$ (se alcançada).
\end{itemize}
A primeira posição que \emph{pode} ser rejeitada é a $3$. Se rejeitada, reamostra-se de $(p-q)_+$: entre os \textit{tokens} listados, apenas ``janela'' tem $p > q$ ($0.70 - 0.12 = 0.58 > 0$; ``porta'' tem $p < q$ e ``parede'' tem $p = q$), de modo que a massa concentra-se em ``janela''. Em seguida, as posições $4$ e $5$ são \textbf{descartadas} (a verificação para na primeira rejeição). Por construção, a saída tem a mesma distribuição do modelo alvo.\\
Pontos-chave: a razão correta é $p/q$ (não $q/p$); estar entre os três mais prováveis do alvo \emph{não} garante aceitação (``porta'' está no \textit{top-3} do alvo, mas $p < q$); após uma rejeição, as posições seguintes são \emph{descartadas}, não reamostradas.\\
\textbf{(a)} confunde estar no \textit{top-3} do alvo com ser aceito (a posição $3$ também pode ser rejeitada).
\textbf{(b)} $p < q$ não implica rejeição certa: há aceitação com prob.\ $0.25$.
\textbf{(c)} inverte a razão ($p/q = 0.25$, não $0.8/0.2$).
\textbf{(e)} $p \ge q$ implica aceitação \emph{garantida} (prob.\ $1$), não rejeição com prob.\ $p/q$.
\textbf{(f)} a reamostragem usa $(p-q)_+$, que exclui ``porta'' ($p < q$); e as posições $4$/$5$ são descartadas, não reamostradas.
\textbf{(g)} se a posição $3$ é rejeitada, a $5$ é descartada; elas não são independentes.
\textbf{(h)} a verificação não exige $p = q$: \textit{tokens} com $p \ge q$ são sempre aceitos.

%% ====================== UNIDADE 5 ======================

\subsection*{Questão 6 \textemdash{} Resposta: (b)}
$\mathrm{Tr}[\boldsymbol{\Sigma}] = 0.5 + 2 = 2.5$; \quad $\boldsymbol{\mu}^{\top}\boldsymbol{\mu} = 2^2 + 0^2 = 4$; \quad $\det\boldsymbol{\Sigma} = 0.5\cdot 2 = 1 \Rightarrow \log\det\boldsymbol{\Sigma} = 0$; \quad $D_z = 2$. Logo
\[
D_{KL} = \tfrac{1}{2}(2.5 + 4 - 2 - 0) = \tfrac{1}{2}(4.5) = 2.25.
\]
\textbf{(a)} $3.25$ esquece o termo $-D_z$.
\textbf{(c)} $4.50$ esquece o fator $\tfrac{1}{2}$.
\textbf{(d)} $0.25$ esquece o termo $\boldsymbol{\mu}^{\top}\boldsymbol{\mu}$.
\textbf{(e)} $1.50$ usa $\det\boldsymbol{\Sigma}$ no lugar de $\mathrm{Tr}[\boldsymbol{\Sigma}]$.
\textbf{(f)} $4.25$ soma $+D_z$ em vez de subtrair.
\textbf{(g)/(h)} demais erros.

\subsection*{Questão 7 \textemdash{} Resposta: (h)}
\textbf{I (V):} decomposição reconstrução $-$ KL; a ELBO é limite \emph{inferior} de $\log p(\mathbf{x})$.
\textbf{II (V):} sem o \textit{reparametrization trick} não se retropropaga através do nó aleatório de amostragem; movê-lo para $\boldsymbol{\epsilon} \sim \mathcal{N}(0,\mathbf{I})$ mantém o caminho por $\boldsymbol{\mu}, \boldsymbol{\Sigma}^{1/2}$ diferenciável.
\textbf{III (V):} $-\log\det\boldsymbol{\Sigma} \to +\infty$ quando $\det\boldsymbol{\Sigma} \to 0$, penalizando variância nula (colapso a massa pontual).
\textbf{IV (F):} no AE determinístico o espaço latente \emph{não} tem distribuição conhecida; amostrar $\mathbf{z}\sim\mathcal{N}(0,\mathbf{I})$ cai em regiões sem suporte e o \textit{decoder} produz lixo. É justamente o que o VAE corrige ao regularizar o posterior.
\textbf{V (F):} a ELBO é limite \emph{inferior} (não superior); a diferença para $\log p(\mathbf{x})$ é $D_{KL}[q(\mathbf{z}|\mathbf{x},\phi)\,\|\,p(\mathbf{z}|\mathbf{x},\thetav)]$ (KL ao posterior \emph{verdadeiro}, não ao \textit{prior}).
Verdadeiras: I, II e III.\\
\textbf{(a)/(c)} incluem IV (falsa).
\textbf{(d)/(e)} incluem V (falsa).
\textbf{(b)} omite III.
\textbf{(f)} troca verdadeiras por falsas.
\textbf{(g)} marca tudo.

\subsection*{Questão 8 \textemdash{} Resposta: (e)}
Com $\exp(\mathbf{s}) = (e^{\ln 2}, e^{0}) = (2, 1)$:
\[
\mathbf{y}^B = (3\cdot 2 + 1,\; 4\cdot 1 + (-1)) = (7,\; 3).
\]
O log-determinante de uma \textit{affine coupling} é $\log\lvert\det J\rvert = \sum_i s_i = \ln 2 + 0 = \ln 2 \approx 0.693$ (a parte copiada contribui com Jacobiano identidade).\\
\textbf{(a)} $\log\lvert\det J\rvert = 2$ reporta $\det J$ (e não $\log\det J$).
\textbf{(b)} $3$ soma $\exp(s_i) = 2 + 1$ em vez de $\sum s_i$.
\textbf{(c)} $\mathbf{y}^B = (6,4)$ esquece o deslocamento $\mathbf{t}$.
\textbf{(d)} $\mathbf{y}^B = (7,4)$ aplica $\mathbf{t}$ só na primeira componente.
\textbf{(f)} não transforma $\mathbf{h}^B$.
\textbf{(g)/(h)} log-determinante incorreto.

\subsection*{Questão 9 \textemdash{} Resposta: (c)}
\textbf{I (V):} pela regra de Bayes com \textit{priors} iguais, $D^{*}(x) = P(x)/(P(x)+P^{*}(x))$.
\textbf{II (V):} substituindo $D^{*}$, a \emph{loss} vira $2\,D_{JS}[P^{*}\|P] - 2\ln 2$, i.e., Jensen--Shannon a menos de constante.
\textbf{III (V):} $\partial(-\ln\sigma)/\partial d = -(1-\sigma) \to -1$ quando $\sigma\to 0$, ao passo que a \emph{loss} original tem gradiente $\to 0$ (satura).
\textbf{IV (V):} suportes disjuntos $\Rightarrow D_{JS} = \ln 2$ (máximo), constante, com gradiente nulo para o gerador.
\textbf{V (F):} é o oposto: o termo de cobertura \emph{satura} (fica $\approx \tfrac{1}{2}P\ln 2$, quase insensível a $P^{*}$), gerando gradiente \emph{fraco} para cobrir modas; essa assimetria é justamente a causa do \textit{mode collapse}.
Verdadeiras: I, II, III e IV.\\
\textbf{(a)/(e)/(f)} omitem alguma verdadeira.
\textbf{(b)} omite II.
\textbf{(d)/(h)} incluem V (falsa).
\textbf{(g)} marca tudo.

\subsection*{Questão 10 \textemdash{} Resposta: (a)}
\[
\tilde{\mathbf{g}}_t = (1 + w)\,\mathbf{g}_t[\mathbf{z}_t, c] - w\,\mathbf{g}_t[\mathbf{z}_t, \varnothing] = (1+3)\cdot 2 - 3\cdot 0.5 = 8 - 1.5 = 6.5.
\]
(Equivalentemente, $\tilde{\mathbf{g}}_t = \mathbf{g}_t[\mathbf{z}_t,c] + w\,(\mathbf{g}_t[\mathbf{z}_t,c] - \mathbf{g}_t[\mathbf{z}_t,\varnothing]) = 2 + 3\cdot 1.5 = 6.5$: extrapola-se a partir da predição \emph{condicionada}.)\\
\textbf{(b)} $5.0$ usa a predição \emph{sem condição} como base: $\mathbf{g}_\varnothing + w(\mathbf{g}_c - \mathbf{g}_\varnothing)$.
\textbf{(c)} $4.0$ inverte os papéis: $w\,\mathbf{g}_c - (1+w)\mathbf{g}_\varnothing$.
\textbf{(d)} $8.0$ esquece o termo $-w\,\mathbf{g}_\varnothing$.
\textbf{(e)} $6.0$ calcula $(1+w)(\mathbf{g}_c - \mathbf{g}_\varnothing)$.
\textbf{(f)} $4.5$ calcula apenas $w\,(\mathbf{g}_c - \mathbf{g}_\varnothing)$.
\textbf{(g)} $2.0$ é a predição condicionada pura ($w = 0$).
\textbf{(h)} $0.5$ é a predição sem condição.

\color{black}

\end{document}
